\section{What can ownership add to regulation and fiscal redistribution?}
\subsection{The unit of comparison is a policy contract}
A legally required grant of new equity changes private property claims. Calling it usownership does not remove its distributive incidence. Conversely, legislation can create property rights, and a tax-financed asset grant changes a household's balance sheet. The useful distinction is between contracts that distribute current income, accumulate enforceable claims, influence production or insure risks. Each can be implemented through legislation, private agreements or combinations of the two.

Consider five comparators. A conventional transfer supplies current spending power and insurance. A capital-income tax shifts part of the financing away from wages. A dividend levy can be earmarked to the very participants who would receive equity. A tax-funded asset grant or subsidy-for-equity arrangement can build legal ownership. Competition, labour and data regulation change rents, bargaining positions and admissible conduct. No single one is an adequate substitute for every function of the others. An ownership requirement does not enforce product safety, and a public income floor does not eliminate monopoly power.

\subsection{A cash-flow equivalence benchmark}
Let $o_{ij}$ denote household $i$'s original ownership fraction of firm $j$, with no subsequent private share trades. Let $U_{j,t}$ be aggregate user equity under \eqref{eq:U}, $u_{ij,t}$ its household allocation, and $\Pi_{j,t}$ the firm's distributable cash. Legacy holdings then equal $(1-U_{j,t})o_{ij}$. The equity system pays
\begin{equation}
D^E_{i,t}=\sum_j\big[(1-U_{j,t})o_{ij}+u_{ij,t}\big]\Pi_{j,t}.\label{eq:equitypayout}
\end{equation}

\begin{proposition}[Matched dividend-levy equivalence]\label{prop:equivalence}
Suppose profits, investment, participant weights, subsequent taxes and information are identical, the levy is enforceable, and shares are not traded. Retain original legal ownership, levy $U_{j,t}$ of firm $j$'s distributable cash pro rata from original holders, and rebate $u_{ij,t}\Pi_{j,t}$ to participant $i$. This system reproduces \eqref{eq:equitypayout} in every period and state. In a model where household decisions depend only on those cash flows and opening cash, both systems generate identical consumption, MVI, money and prices.
\end{proposition}
\begin{proof}
Original-holder income after the levy is $\sum_j(1-U_{j,t})o_{ij}\Pi_{j,t}$. Adding the rebate gives \eqref{eq:equitypayout}. Given equal initial cash, the next cash balance and all income-dependent decisions coincide. Induction establishes the path result. This argument holds conditional on identical investment; it does not assume that a real tax and a real equity contract necessarily induce identical behaviour.
\end{proof}

The proposition is an intentionally demanding null benchmark. The simulations implement the levy through a \emph{shadow entitlement register}, using the same contribution history, dilution schedule and investment wedge as usownership. The gross earmarked levy equals its rebate; it is reported separately from conventional tax revenue and MVI. It is not fictitious additional seigniorage. The comparator is sophisticated rather than a description of a current tax system. Its purpose is to show what an equity label alone cannot establish.

The same result qualifies the stock-versus-flow argument. If identical future tax liabilities and benefit claims are enforceable and capitalised at the same state-contingent discount factors, economic net wealth is also identical. A lower Gini of \emph{legal book assets}, which omits those future fiscal claims, is not by itself a welfare gain. Legal ownership can still matter for votes, sale, collateral, inheritance and protection against policy revision. These are additional contractual rights, not new real output created by distributing certificates.

\subsection{Gradual does not mean inexpensive}
At fixed firm value, a grant of $\delta$ new shares per old share transfers $\delta/(1+\delta)$ of the post-issue claim. At constant annual issuance of 0.5\%, 1\% or 2\%, user shares after 50 years are respectively 22.1\%, 39.2\% and 62.8\%. The small yearly rate is politically and operationally different from an immediate large transfer, but its compound incidence is not small.

For a transparent valuation calculation, let total end-year dividends grow at $g_D$ and let investors discount at $r>g_D$. The ratio of legacy-share present value with perpetual constant dilution to its no-dilution value is
\begin{equation}
\frac{V_\delta}{V_0}=\frac{(1+r)-(1+g_D)}{(1+r)(1+\delta)-(1+g_D)}.\label{eq:pv}
\end{equation}
This follows by summing two geometric series with ratios $(1+g_D)/(1+r)$ and $(1+g_D)/[(1+r)(1+\delta)]$. With the illustrative assumptions $r=8\%$ and $g_D=3\%$, the legacy-value losses are 9.7\%, 17.8\% and 30.2\% at the three issuance rates. These are mechanical perpetuity values, not estimated financing-cost effects. The same burden arises under a fully anticipated matched levy.

At a particular horizon, retaining the old shareholders' terminal wealth requires total company value to rise by $1/(1-U_T)-1$ relative to the no-grant counterfactual. Thus 1\% annual issuance over 50 years would require about 64.5\% additional total value to offset its ownership dilution. It would be unjustified to assume that stronger engagement automatically supplies this gain. Supplementary Table S1 reports the complete incidence arithmetic.

\subsection{Conditional advantages worth testing}
Four margins are promising. \emph{First, persistent claims}: a vested share may survive the expiry of a programme's new-award budget. A durable entitlement can reduce renegotiation and help participants plan long-term investment in an ecosystem. A legally protected tax-funded endowment may do the same; corporate rights are themselves vulnerable to insolvency, recapitalisation and changes in law. The relative security of the two contracts must be measured.

\emph{Second, participation and governance}: a participant who contributes expertise, useful feedback or productive complements can share a firm's residual upside and, where rights permit, its governance. That may reward contributions inadequately covered by conventional employment or purchase contracts. Fixed participation scores in our macro model do not estimate the resulting effort response. Non-voting shares should not be presented as democratic control, and fragmented minority positions may require representation to exercise any practical influence.

\emph{Third, issuance-date liquidity}: a new-share grant need not require a firm or government to pay cash at the issue date. Compared with a cash-financed share purchase, it can be useful where liquidity is scarce even though its economic cost falls on existing holders. It can also attach claims before profits are distributed. Whether a cash-flow tax with loss offsets, deferred payments or equity settlement is equally effective remains a counterfactual, not an impossibility.

\emph{Fourth, a verifiable accumulating claim}: a disclosed share fraction and register can make the allocation rule inspectable and avoid repeatedly determining an unlisted firm's market value merely to compute the number of shares. That does not remove valuation from accounting, tax or disputes. Contribution verification, group consolidation, corporate actions, buybacks and anti-dilution protection create their own administrative burdens. No percentage-of-GDP bureaucracy saving is assumed.

A useful conditional test is whether the present value of additional participant value, governance gains, commitment and financing relief exceeds measurement costs, concentrated risk and investment distortions. Denote those components by $V_p,V_g,V_c,V_l$ and $C_m,C_r,C_i$. The proposal's institutional case requires
\begin{equation}
\Delta W_{\rm institutional}=V_p+V_g+V_c+V_l-C_m-C_r-C_i>0.\label{eq:wedge}
\end{equation}
Equation~\eqref{eq:wedge} organises quantities to estimate; it is not an estimated welfare equation or an assumption that the left side is positive. The frictionless replication benchmark sets the differential terms to zero. This is a stronger empirical research agenda than treating all conventional policy as incapable of accumulating assets.

\subsection{Why the Internet-of-Value paradigm changes the practical comparison}
The cash-flow equivalence proposition is deliberately technology-neutral. In a frictionless model, legal equity and a perfectly credible matched levy can deliver identical cash. The practical comparison changes when the institutional cost of maintaining those claims differs. The existing financial architecture is highly capable for large conventional securities, but repeated low-value allocations to millions of heterogeneous users are not its standard use case. A tokenised infrastructure can represent fractional claims, transfer conditions, vesting and payment instructions in the same programmable environment \citep{bis2023,bis2025}.

This produces a conditional implementation advantage rather than an economic free lunch. Under usownership, a user's verified contribution can update an accumulating property claim without the state repeatedly collecting cash and purchasing assets on the user's behalf. Dividends can be routed to the current owner, corporate actions can update the same register, and an eventual voluntary exchange into diversified assets can settle against tokenised money. Under a conventional levy, the fiscal authority can reproduce income but must maintain a separate entitlement rule and political commitment. Either architecture can be designed badly; programmability simply lowers some contracting and reconciliation frictions when assets themselves are digital and executable.

The structural-change argument is therefore stronger than a blockchain argument. As economic value migrates from documents and disconnected databases toward digitally native or tokenised claims, ownership design becomes part of the software and legal architecture through which value moves. The European Commission now describes tokenisation as a possible path toward an Internet of Value, while the BIS sees programmable money and assets as components of a next-generation monetary system \citep{eciov2026,bis2025}. Usownership is tailored to that environment because the distribution mechanism is itself a stream of small ownership events. Its comparative advantage should disappear in the model if programmable infrastructure does not reduce implementation costs or if equivalent fiscal entitlements provide the same durability and governance.

The technology also creates new failure modes. An erroneous contribution oracle can automate the wrong allocation at scale; a compromised key can transfer a legally valuable claim; privacy-invasive scoring can turn ownership into surveillance; and incompatible ledgers can recreate fragmentation. These risks belong inside the cost terms of Equation~\eqref{eq:implementation}. A public permissionless blockchain is neither necessary nor assumed.

\subsection{Monetary finance is an auxiliary bridge, not a unique bridge}
Declining labour's share can raise distributional pressures without reducing the absolute real wage bill: sufficiently rapid output growth can offset the share decline. The Commission's 2026 tax-trends update reports 2024 EU taxes at 39.4\% of GDP, of which labour taxes and social contributions provide 51.5\% \citep{taxtrends2026}. Their product is 20.29\% of GDP. At fixed GDP, rates and composition, a 25\% or 50\% contraction in this base would reduce receipts by 5.07 or 10.15 percentage points of GDP. These are static exposures, not forecasts; neither job exposure nor a falling labour share implies either contraction.

A government can bridge an income disruption with redesigned taxes, debt, existing reserves, public asset income or seigniorage. Borrowing may be especially relevant to a temporary shock. Because the laboratory excludes sovereign debt, it cannot establish a timing advantage of monetary finance over debt-financed transfers. The useful monetary question is more limited: how much financing can coexist with sustained transaction demand and acceptable prices, after accounting for the incidence of any holding levy? The monetary accounting in the following section gives a bounded answer rather than an instruction to replace fiscal policy with a programmable ledger.
